% Auto-generated by sa_generator.py
% Compile with: lualatex sa_output.tex
\documentclass[11pt,letterpaper,openany]{scrbook}

\usepackage[letterpaper,
  bindingoffset=0.2in,
  left=0.45in, right=0.45in,
  top=0.4in, bottom=0.4in,
  footskip=0.2in, marginparwidth=0em
]{geometry}

\usepackage{fontspec}
\usepackage{polyglossia}
\usepackage{xcolor}
\usepackage{ragged2e}
\usepackage{titlesec}
\usepackage{microtype}

\setmainlanguage{english}
\setotherlanguage{hebrew}

% ── Colors ──
\definecolor{hcolor}{HTML}{1A3264}
\definecolor{ruleGrey}{HTML}{AAAAAA}
\definecolor{lightGrey}{HTML}{E0E0E0}

% ── Fonts ──
\newfontfamily\hebrewfont[
  Path=../Fonts/, Script=Hebrew, Ligatures=TeX,
  BoldFont=FrankRuehlCLM-Bold.otf
]{FrankRuehlCLM-Medium.otf}

\newfontfamily\mechaberfont[
  Path=../Fonts/Frank/static/, Script=Hebrew, Ligatures=TeX,
  BoldFont=FrankRuhlLibre-Bold.ttf
]{FrankRuhlLibre-Regular.ttf}

\newfontfamily\commentaryfont[
  Path=../Fonts/, Script=Hebrew, Ligatures=TeX,
  BoldFont=FrankRuehlCLM-Bold.otf
]{FrankRuehlCLM-Medium.otf}

\newfontfamily\engfont[
  Path=../Fonts/, Ligatures=TeX,
  BoldFont=EBGaramond-Bold.ttf,
  ItalicFont=EBGaramond-Italic.ttf
]{EBGaramond-Regular.ttf}

% ── Spacing ──
\setlength\parindent{0pt}
\setlength{\columnsep}{1em}
\raggedbottom
\pagestyle{empty}

% ── Interlinear macros ──
% Hebrew catchword followed by em-dash, then translation follows inline
\newcommand{\Heb}[1]{\texthebrew{\textbf{#1}}~--- }
% Seif marker: just the Hebrew letter, large and bold
\newcommand{\Seif}[1]{{\Large\bfseries\texthebrew{#1}}\ }
% Rema marker
\newcommand{\Rema}{\texthebrew{\textbf{הגה:}}\ }
% Commentary header
\newcommand{\CommHeader}[1]{%
  \par\noindent{\commentaryfont\fontsize{9}{11}\selectfont%
  \textcolor{hcolor}{\texthebrew{\textbf{#1}}}}\par\smallskip
}
% Seif katan marker in commentary
\newcommand{\SK}[1]{{\bfseries\texthebrew{#1}}\ }
% Thin rule
\newcommand{\thinrule}{\par\noindent\textcolor{ruleGrey}{\rule{\linewidth}{0.3pt}}\par\vspace{2pt}}

\begin{document}
\begin{sloppypar}

%% ━━━━ PAGE 1 ━━━━
%% Header
\noindent
\colorbox{lightGrey}{\parbox{\dimexpr\textwidth-8pt\relax}{%
  \centering\vspace{3pt}%
  \begin{hebrew}%
  {\mechaberfont\fontsize{11}{13}\selectfont%
  \textcolor{hcolor}{%

שולחן ערוך אורח חיים \quad הלכות שבת \quad סימן ר"צ \quad דף קנ"ט
}}%
  \end{hebrew}%
  \vspace{3pt}%
}}

\vspace{6pt}


%% 3-column layout: Left comm | Center Mechaber | Right comm
\noindent
%% LEFT: Magen Avraham
\begin{minipage}[t]{0.235\textwidth}
\raggedright

\CommHeader{מגן אברהם}
{\commentaryfont\fontsize{8.5}{11}\selectfont
\Heb{אין להניח הדרשה עבור סעודה ג' כדאי' במד' משלי בפ' אשת חיל שמתו בניו של ר"מ מאותו החטא [רוקח]}
}
{\engfont\fontsize{8}{10}\selectfont One should not forgo the public Torah discourse in favor of the third meal, as we find in the Midrash on Mishlei, Parshas Eishes Chayil, that the sons of Rabbi Meir died because of this sin [Rokeach].}
\par\smallskip
{\commentaryfont\fontsize{8.5}{11}\selectfont
\Heb{וצ"ל שהדרשה תהא להורות לעם את חקי האלהים ואת תורותיו ולהכניס יראת שמים בלבבם ולא כמו שנוהגין עכשיו}
}
{\engfont\fontsize{8}{10}\selectfont And one must say that the discourse should be to instruct the people in the statutes of God and His laws, and to instill the fear of Heaven in their hearts --- and not as is the practice nowadays.}
\par\smallskip
{\commentaryfont\fontsize{8.5}{11}\selectfont
\Heb{לא יאמר בשבת נישן כדי לעשות מלאכתינו בערב [ס"ח סימן רס"ו]}
}
{\engfont\fontsize{8}{10}\selectfont One should not say on Shabbos, 'Let us sleep so that we can do our work in the evening' [Sefer Chassidim, siman 266].}
\par\smallskip
{\commentaryfont\fontsize{8.5}{11}\selectfont
}

\end{minipage}%
\hfill
%% CENTER: Mechaber + Rema
\begin{minipage}[t]{0.47\textwidth}

{\mechaberfont\fontsize{12}{15.5}\selectfont
\Seif{א}
\Heb{ירבה פירות ומגדים ומיני ריח כדי להשלים מנין מאה ברכות:}
}
{\engfont\fontsize{10}{13}\selectfont One should increase [consumption of] fruits, delicacies, and fragrant spices in order to complete the count of one hundred blessings.}
\par\medskip
{\mechaberfont\fontsize{12}{15.5}\selectfont
\Rema
\Heb{ואם רגיל בשינת צהרים אל יבטלנו כי עונג הוא לו (טור):}
}
{\engfont\fontsize{10}{13}\selectfont And if one is accustomed to an afternoon nap, he should not forgo it, for it is a delight for him (Tur).}
\par\medskip
{\mechaberfont\fontsize{12}{15.5}\selectfont
\Seif{ב}
\Heb{אחר סעודת שחרית קובעים מדרש לקרות בנביאים ולדרוש בדברי אגדה ואסור לקבוע סעודה באותה שעה:}
}
{\engfont\fontsize{10}{13}\selectfont After the morning meal, one establishes a study session to read from the Prophets and to expound on Aggadic teachings. It is forbidden to schedule a meal during that time.}
\par\medskip
{\mechaberfont\fontsize{12}{15.5}\selectfont
\Rema
\Heb{ופועלים ובעלי בתים שאינן עוסקים בתורה כל ימי שבוע יעסקו יותר בתורה בשבת מתלמידי חכמים העוסקים בתורה כל ימי השבוע ותלמידי חכמים ימשיכו יותר בעונג אכילה ושתיה קצת שהרי מתענגים בלמודם כל ימי השבוע [ב"י סי' רפ"ח בשם ירושלמי]:}
}
{\engfont\fontsize{10}{13}\selectfont Workers and householders who do not engage in Torah study during the weekdays should study more Torah on Shabbos than Torah scholars who study all week. And Torah scholars should indulge somewhat more in the pleasure of eating and drinking, since they already derive pleasure from their studies throughout the week [Beis Yosef, Siman 288, citing the Yerushalmi].}
\par\medskip

\end{minipage}%
\hfill
%% RIGHT: Taz / Magen David
\begin{minipage}[t]{0.235\textwidth}
\raggedright

\CommHeader{מגן דוד}
{\commentaryfont\fontsize{8.5}{11}\selectfont
\Heb{ירבה פירות ומגדים. הנה ידוע כי מספר מאה ברכות הוא תקנת דוד המלך עליו השלום כמבואר בגמרא מנחות דף מ"ג ע"ב}
}
{\engfont\fontsize{8}{10}\selectfont One should increase fruits and delicacies. It is well known that the count of one hundred blessings was enacted by King David, peace be upon him, as explained in the Gemara, Menachos 43b.}
\par\smallskip
{\commentaryfont\fontsize{8.5}{11}\selectfont
\Heb{ובשבת שחסרים כמה ברכות מתפילת שמונה עשרה שמתפללים רק שבע ברכות ישלים אותם ע"י ברכות הנהנין על פירות ובשמים ומיני מגדים}
}
{\engfont\fontsize{8}{10}\selectfont And on Shabbos, when several blessings are missing from the Shemoneh Esrei --- for one recites only seven blessings --- one should complete them through blessings of enjoyment over fruits, fragrances, and delicacies.}
\par\smallskip
{\commentaryfont\fontsize{8.5}{11}\selectfont
\Heb{ויכוין בכל ברכה וברכה שמברך לשם שמים ולהשלמת המאה ברכות שתיקן דוד המלך}
}
{\engfont\fontsize{8}{10}\selectfont And one should have intent with each and every blessing that he recites, for the sake of Heaven and for the completion of the one hundred blessings that King David enacted.}
\par\smallskip
{\commentaryfont\fontsize{8.5}{11}\selectfont
}

\end{minipage}

\newpage
\begin{center}
{\engfont\Large\textcolor{hcolor}{Page 2 --- Commentary Translations (TODO)}}
\end{center}

\end{sloppypar}
\end{document}
